\chapter{Diseño de Base de Datos --- Caso Práctico}

\section{Proceso de diseño}

\begin{enumerate}
  \item Recopilar requisitos (entrevistas con stakeholders)
  \item Modelo conceptual (ER / UML)
  \item Modelo lógico (tablas, atributos, tipos, FK)
  \item Normalización (3NF mínimo)
  \item Modelo físico (índices, particiones, tipos específicos del DBMS)
  \item Revisión de rendimiento (EXPLAIN, pruebas de carga)
\end{enumerate}

\section{Caso: Sistema de Biblioteca}

\textbf{Requisitos:} registrar libros, copias físicas, socios, préstamos y devoluciones.

Entidades: \texttt{Libro}, \texttt{Copia}, \texttt{Socio}, \texttt{Prestamo}, \texttt{Autor}, \texttt{Categoria}.

Relaciones clave:
\begin{itemize}
  \item Libro tiene muchas Copias (1:N)
  \item Libro tiene muchos Autores (M:N via \texttt{libro\_autor})
  \item Copia puede estar en muchos Préstamos a lo largo del tiempo (1:N)
  \item Socio tiene muchos Préstamos (1:N)
\end{itemize}

\begin{diagramabox}{ER --- Sistema Biblioteca}
\begin{center}
  % Generar con: make diagrams (requiere Java + plantuml.jar)
  % Fuente: images/src_plantuml/11_er_biblioteca.puml
  \fbox{\parbox{0.7\textwidth}{\centering\vspace{1em}[Diagrama PlantUML: ejecutar \texttt{make diagrams}]\vspace{1em}}}
\end{center}
\end{diagramabox}

\begin{lstlisting}[caption={DDL completo --- Sistema Biblioteca}]
CREATE TABLE categorias (
    id     SERIAL PRIMARY KEY,
    nombre VARCHAR(80) NOT NULL UNIQUE
);

CREATE TABLE libros (
    id           SERIAL PRIMARY KEY,
    isbn         VARCHAR(13) UNIQUE NOT NULL,
    titulo       VARCHAR(200) NOT NULL,
    anio_pub     SMALLINT,
    id_categoria INTEGER REFERENCES categorias(id)
);

CREATE TABLE autores (
    id     SERIAL PRIMARY KEY,
    nombre VARCHAR(150) NOT NULL
);

CREATE TABLE libro_autor (
    id_libro  INTEGER REFERENCES libros(id)  ON DELETE CASCADE,
    id_autor  INTEGER REFERENCES autores(id) ON DELETE CASCADE,
    PRIMARY KEY (id_libro, id_autor)
);

CREATE TABLE copias (
    id        SERIAL PRIMARY KEY,
    id_libro  INTEGER NOT NULL REFERENCES libros(id),
    codigo    VARCHAR(20) UNIQUE NOT NULL,
    estado    VARCHAR(20) DEFAULT 'disponible'
              CHECK (estado IN ('disponible','prestada','dannada','baja'))
);

CREATE TABLE socios (
    id        SERIAL PRIMARY KEY,
    nombre    VARCHAR(150) NOT NULL,
    email     VARCHAR(200) UNIQUE NOT NULL,
    activo    BOOLEAN DEFAULT TRUE,
    creado_en DATE DEFAULT CURRENT_DATE
);

CREATE TABLE prestamos (
    id          SERIAL PRIMARY KEY,
    id_copia    INTEGER NOT NULL REFERENCES copias(id),
    id_socio    INTEGER NOT NULL REFERENCES socios(id),
    fecha_ini   DATE NOT NULL DEFAULT CURRENT_DATE,
    fecha_dev   DATE NOT NULL,
    devuelto_en DATE,
    CONSTRAINT chk_fechas CHECK (fecha_dev > fecha_ini)
);
\end{lstlisting}

\begin{lstlisting}[caption={Consultas operacionales sobre el diseño}]
-- Copias disponibles por categoría
SELECT c.nombre AS categoria, COUNT(cp.id) AS copias_disponibles
FROM   categorias c
JOIN   libros l  ON l.id_categoria = c.id
JOIN   copias cp ON cp.id_libro = l.id AND cp.estado = 'disponible'
GROUP BY c.nombre
ORDER BY copias_disponibles DESC;

-- Socios con préstamos vencidos
SELECT s.nombre, s.email, p.fecha_dev, l.titulo
FROM   prestamos p
JOIN   socios s  ON p.id_socio = s.id
JOIN   copias cp ON p.id_copia = cp.id
JOIN   libros l  ON cp.id_libro = l.id
WHERE  p.devuelto_en IS NULL
  AND  p.fecha_dev < CURRENT_DATE
ORDER BY p.fecha_dev;
\end{lstlisting}

\begin{entrevistabox}
\begin{enumerate}
  \item ¿Cómo modelarías la diferencia entre un libro (título) y una copia física en una biblioteca?
  \item ¿Qué índices agregarías a este diseño para optimizar las consultas más frecuentes?
  \item ¿Cómo manejarías el historial de préstamos si un socio puede tomar el mismo libro varias veces?
\end{enumerate}
\end{entrevistabox}
